\section{Introduction}\label{sec:introduction}

Let $(X,\mathcal F,\mu,T)$ be an invertible probability-preserving system.
For bounded measurable functions $f_1,\ldots,f_d$, the multiple ergodic
averages are
\[
A_N(f_1,\ldots,f_d)(x)
=\frac1N\sum_{n=1}^N\prod_{j=1}^d f_j(T^{jn}x).
\]
The general pointwise question asks whether these averages converge
almost everywhere for every finite $d$ on every invertible
probability-preserving system. Here we treat four functions under
the mixing hypothesis. We say that $T$ is mixing if
\begin{equation}\label{eq:mixing}
 \mu(A\cap T^{-n}B)\longrightarrow\mu(A)\mu(B)
 \qquad (|n|\longrightarrow\infty)
\end{equation}
for every $A,B\in\mathcal F$. The functions are fixed throughout each
convergence assertion. No rate in \eqref{eq:mixing} is assumed.

\begin{theorem}\label{thm:main}
Let $(X,\mathcal F,\mu)$ be a probability space and let $T:X\to X$
be invertible, bimeasurable, measure preserving, and mixing in the sense
of \eqref{eq:mixing}. For every $f_1,f_2,f_3,f_4\in L^\infty(\mu)$,
\[
 \frac1N\sum_{n=1}^N
 f_1(T^nx)f_2(T^{2n}x)f_3(T^{3n}x)f_4(T^{4n}x)
 \longrightarrow \prod_{j=1}^4\int_X f_j\dd\mu
\]
for $\mu$-almost every $x$, along all positive integers $N$.
The original probability space need not be standard.
\end{theorem}

\subsection*{Context and inputs}

Birkhoff's pointwise ergodic theorem treats a single observation along
an orbit \cite{Birkhoff}. Furstenberg's ergodic proof of Szemer\'edi's
theorem made products at arithmetic-progression times a central object
of study \cite{Furstenberg}. For weakly mixing transformations, these
averages converge in $L^2$ to the product of the integrals. For arbitrary
transformations, Host and Kra established norm convergence of every
finite consecutive length using nilpotent characteristic factors
\cite[Theorem~1.1]{HostKra}; Ziegler independently developed universal
characteristic factors and another proof \cite{Ziegler}. Norm
convergence alone does not control the exceptional lengths at an
individual point.

Bourgain proved pointwise convergence for two powers of one
transformation \cite{Bourgain}. Higher-order pointwise results include
$K$ systems \cite[Theorem~4.2]{DerrienLesigne}, weakly mixing systems
with suitable spectral restrictions \cite{Assani}, and weakly mixing
systems whose pairwise-independent finite self-joinings are all
product measures \cite{GutmanHuangShaoYe}. Huang, Shao and Ye proved
pointwise convergence of every finite consecutive length on ergodic
measure-distal systems \cite[Theorem~C]{HuangShaoYe}, extending the
distal investigation of Lesigne \cite{Lesigne}. Recent pointwise
theorems for polynomial iterates include the bilinear result of
Krause, Mirek and Tao \cite{KrauseMirekTao} and the distinct-degree
multilinear result of Kosz, Mirek, Peluse, Wan and Wright
\cite{KoszMirekPeluseWanWright}. Their degree hypotheses do not cover
the four linear iterates considered here. Kuca's survey \cite{Kuca}
gives a broader account of the norm and pointwise questions.

Theorem~\ref{thm:main} uses three companion results. The multiple-mixing
theorem of \cite[Theorem~1.1]{MixingCompanion} supplies mixing of every finite order
from \eqref{eq:mixing}. The analytic construction of
\cite{TripleCompanion}, based on the local estimates of
\cite{HilbertCompanion}, supplies an oscillation estimate for threefold
averages. Its oscillation assertion holds on arbitrary invertible
systems; mixing is used separately to identify limits. We need the
threefold estimate at each three-element subset of $\{1,2,3,4\}$.
Appendix~\ref{sec:lower-length} proves the required coefficient
adaptation and the resulting pointwise convergence lemma, including its
nonergodic formulation. The local analytic theorems themselves are cited
with their hypotheses and locations. Thus the contribution here is the
fourfold deduction from these specified inputs.

The structural part uses the Furstenberg--Zimmer theory of relatively
compact and relatively weakly mixing extensions. We record the
nonergodic and power-stability statements at their point of use in
Section~\ref{sec:structure}. The construction also uses extensions on
which certain conditional expectations cannot improve. Such energy
maximization belongs to the same general method as sated extensions
\cite{Austin}; the data being saturated here come from selected
four-coordinate laws.

\subsection*{The difficulty caused by selected lengths}

Suppose that the conclusion fails. After subtracting means, there are
four centred real functions and a positive-measure set of points at
which their averages remain positively separated from zero along
arbitrarily late lengths. We choose one such length $N_a(x)$ after
scale $2^a$. Averaging the configurations
\[
 (T^nx,T^{2n}x,T^{3n}x,T^{4n}x),\qquad 1\le n\le N_a(x),
\]
and passing to a limit preserves a nonzero correlation. A constant-roof
suspension makes these limit laws invariant under a flow and permits
the changes of speed needed later.

Selection is the main analytic obstruction. Conditioning on a length
chosen from the initial state can change the joint law of independent
auxiliary variables. Ordinary norm convergence, with no
rate, does not by itself bound all the resulting errors. We introduce
three groups of auxiliary spatial coordinates. Their joint law is a
product, and any two groups together are independent of the original
coordinate, after conditioning on the suspension phases. A length
event may change the full three-group law, but preserves the
two-group marginals. This is precisely the amount of independence used
in Section~\ref{sec:channels}. For a finite model of the distinction,
let $U,V,W$ be independent symmetric signs. Their product is independent
of each pair, but conditioning on $UVW=1$ constrains the triple.

The estimate making this possible is a finite-dimensional packing
statement proved in Section~\ref{sec:packing}. If two bounded families
of Hilbert-space vectors have small pairwise inner products outside a
graph of bounded degree, few of their tensor products can have large
projection onto a subspace of slowly growing dimension. A Gaussian
comparison and a trace-moment estimate give a quantitative form of
this assertion. Applied to the leading covariance subspace, it controls
the projected part of each selected average. The complementary part
has a square-summable error over the possible lengths. The estimate is
independent of the ambient Hilbert-space dimensions and requires no
quantitative mixing rate.

\subsection*{From a correlation to a contradiction}

Each auxiliary system carries both an original spatial coordinate and
the three groups of additional coordinates. We call such a system a
channel after specifying its exact independence conditions in
Definition~\ref{def:channel}. In its selected four-coordinate law,
consider nonzero correlations in which each slot uses either a
centred function of the original coordinate or a centred function of
the additional coordinates. Bounded functions of the distal parameters
are allowed as weights. The initial failure gives a correlation with
no additional-coordinate slots.

Choose a channel and a correlation maximizing the number of slots of
the second kind. Independence excludes four such slots. An iterated
array construction excludes three: products along different axes
through a fixed corner are orthogonal, but each has the same nonzero
correlation with the function at that corner. Bessel's inequality
rules this out. At most two slots can therefore be of the second kind.

The final construction increases that number. First enlarge the
channel until the conditional expectations needed by the argument
are stable under further extensions. Then regard one selected
four-coordinate law as a new system and select once more. The result
is a square of old states. Its distinguished row and column meet at
one corner; saturation supplies their exact conditional laws. A
reanchoring identity on the distal parameters shows that the row and
column correlation cannot vanish. Regrouping the other coordinates
therefore replaces one original-coordinate slot by an
additional-coordinate slot, contradicting maximality.

Sections~\ref{sec:structure} and~\ref{sec:selected} establish the
structural inputs and selected laws. The packing estimate and channel
preservation are proved in Sections~\ref{sec:packing}
and~\ref{sec:channels}. Section~\ref{sec:activation} carries out the
maximality argument and completes the proof. The coefficient-dependent
analytic work is collected in Appendix~\ref{sec:lower-length}.
